\documentclass[a4paper,12pt]{article} 

%%%%%%%%%%%%%%%%%%%%%%%%%%%%%%%% CONSTANTES %%%%%%%%%%%%%%%%%%%%%%%%%%%%%%%%%%%
\newcommand{\numero}{1}                                    %Numéro de la série -1

\usepackage[french]{babel}
\usepackage[utf8]{inputenc}
\usepackage{answers}

\usepackage{hyperref}
\usepackage{multicol}

\usepackage[table,xcdraw]{xcolor}
\usepackage{listings}
\definecolor{ForestGreen}{RGB}{34,139,34}


\usepackage{enumitem}

\AtBeginDocument{\def\labelitemi{$\bullet$}}


\newcommand{\py}{\lstinline{Python} }


\definecolor{backcolour}{rgb}{0.95,0.95,0.92}

\lstset{%
	language         = Python,
	backgroundcolor  = \color{backcolour},
	basicstyle       = \ttfamily, % \upshape\ttfamily,
	keywordstyle     = \bfseries\color{blue}, %\bfseries,
	stringstyle      = \color{magenta},
	commentstyle     = \color{ForestGreen},
	alsoletter = > ,
	morekeywords = {>>>,as,assert,False,None, nonlocal,True, with,yield , <<, >>, :},
	showstringspaces = false,
	numbers=left,
	stepnumber=1,
	literate={à}{{\`{a}}}1 {é}{{\'e}}1 {è}{{\`{e}}}1 {ê}{{\^{e}}}1 {Ê}{{\^{E}}}1 {î}{{\^i}}1 {ô}{{\^{o}}}1 {ç}{{\c{c}}}1 {Ç}{{\c{C}}}1
}

\newcommand{\itemb}[1]{\item \textbf{#1}}

\usepackage{fancyhdr}  %package pour en-tetes et pied de pages
\usepackage{sectsty} % Permet de faire des modifications de police dans diverses sections des "headings" (cf. modif presentation de la page)
\pagestyle{fancy}       %Style pour en-tetes et pieds de pages
\fancyhead[CO,CE]{\sc Série 1\hspace{0.5mm}}
\fancyhead[RO,LE]{Collège Sismondi}  % LaTeX/TEX define \strut to be an invisible box of width zero that extends just enough above and below the baseline. Cela permet d'augementer légèrement la taille en bas de la box de manière à ce qu'elle soit collée à la ligne.
\fancyhead[LO,RE]{\small\ \textsl{1\textsuperscript{ère} année - DO Informatique}}
\fancyfoot[RO,LE]{2021 - 2022}
\fancyfoot[LO,RE]{\small }
\fancyfoot[CO,CE]{\thepage}

\fancyhfoffset[l]{1.2cm} % le "l" en paramètre permet d'indiquer qu'on ne veut modifier que la marge à gauche.
\renewcommand{\headrule}{{%
		\hrule \headwidth \headrulewidth \vskip-\headrulewidth}}
\renewcommand\footrulewidth{\headrulewidth}
\renewcommand{\footrule}{{%
		\vskip-\footruleskip\vskip-\footrulewidth
		\hrule \headwidth \footrulewidth\vskip\footruleskip}}

\usepackage{tikz}
%-------------------------------------------------------------------------------
%---- Eclairage : en encadré sur fond jaune avec symbôle "ampoule" à gauche ----
%-------------------------------------------------------------------------------
\definecolor{coleclairage}{RGB}{255 , 221 , 156}
\definecolor{contoureclairage}{RGB}{255 , 192 , 0}
\newenvironment{eclairage}
{
	\begin{center}%
		\begin{tikzpicture}%
			\node[rectangle, draw=contoureclairage, top color=coleclairage!50, bottom color=coleclairage!140, rounded corners=5pt, inner xsep=5pt, inner ysep=6pt, outer ysep=10pt]\bgroup                     
			\begin{minipage}{0.98\linewidth}
				\begin{minipage}{0.08\linewidth}\centerline{\includegraphics[scale=1]{Symbole_eclairage.png}}\end{minipage}
				\begin{minipage}{0.89\linewidth}\itshape\footnotesize
				}
				{                		
				\end{minipage}
			\end{minipage}\egroup;%
		\end{tikzpicture}%
	\end{center}%
}

%-------------------------------------------------------------------------------
%---- apprendre : en encadré sur fond jaune avec symbôle "ampoule" à gauche ----
%-------------------------------------------------------------------------------
\definecolor{colapprendre}{RGB}{50,205,50}
\definecolor{contourapprendre}{RGB}{34,139,34}
\newenvironment{apprendre}
{
	\begin{center}%
		\begin{tikzpicture}%
			\node[rectangle, draw=contourapprendre, top color=colapprendre!10, bottom color=colapprendre!50, rounded corners=5pt, inner xsep=5pt, inner ysep=6pt, outer ysep=10pt]\bgroup                     
			\begin{minipage}{0.98\linewidth}
				\begin{minipage}{0.08\linewidth}\centerline{\includegraphics[width=30px]{Symbole_learn.png}}\end{minipage}
				\begin{minipage}{0.89\linewidth}\itshape\footnotesize
				}
				{                		
				\end{minipage}
			\end{minipage}\egroup;%
		\end{tikzpicture}%
	\end{center}%
}

\definecolor{colimportant}{RGB}{247 , 189 , 164}
\definecolor{contourimportant}{RGB}{237 , 125 , 49}
\newenvironment{important}
{
	\begin{center}%
		\begin{tikzpicture}%
			\node[rectangle, draw=contourimportant, top color=colimportant!50, bottom color=colimportant!140, rounded corners=5pt, inner xsep=5pt, inner ysep=6pt, outer ysep=10pt]\bgroup                     
			\begin{minipage}{0.08\linewidth}\centerline{\includegraphics[scale=0.8]{Symbole_attention.png}}\end{minipage}
			\begin{minipage}{0.89\linewidth}
			}
			{                		
			\end{minipage}\egroup;
		\end{tikzpicture}%
	\end{center}%
}

%-----------------------------------------------------------------
%---- Modification présentation de la page: marges de la page ----
%-----------------------------------------------------------------
%\addtolength{\hoffset}{-1in}              % 1
%\addtolength{\voffset}{-1in}              % 2
\addtolength{\oddsidemargin}{-0.1 in} % 3
\addtolength{\evensidemargin}{-1in} % 3
\addtolength{\topmargin}{-1in}       % 4
\addtolength{\headheight}{6pt}       % 5
%\addtolength{\headsep}{-0.2cm}           % 6
\setlength{\textheight}{26cm}    % 7
\setlength{\textwidth}{16.5cm}      % 8
\addtolength{\marginparsep}{0pt}      % 9
\setlength{\marginparwidth}{0pt}   % 10
\addtolength{\footskip}{-1mm}           %11

\setlength{\parindent}{0em}% pas d'indentation


% Customiser le nom des sections
\usepackage{titlesec}
\titleformat{\section}[hang]{\Large \bfseries}{Série \thesection:\ }{0pt}{}

\renewcommand{\familydefault}{\sfdefault} % pour avoir des polices san serif

\newtheorem{Exc}{Exercice}
\Newassociation{correction}{Soln}{mycor}
\renewcommand{\Solnlabel}[1]{\bfseries Ex #1 }
\def\exo#1{%
	\futurelet\testchar\MaybeOptArgmyexoo}
\def\MaybeOptArgmyexoo{
	\ifx[\testchar \let\next\OptArgmyexoo
	\else \let\next\NoOptArgmyexoo \fi \next}
\def\OptArgmyexoo[#1]{%
	\begin{Exc}[#1]\normalfont}
	\def\NoOptArgmyexoo{%
		\begin{Exc}\normalfont}
		\newcommand{\finexo}{\end{Exc} \vspace{3mm}}
	\newcommand{\flag}[1]{}
	\newcommand{\entete}[1]

\newcommand{\getexocompteur}{{\the\numexpr \arabic{Exc}  \relax}}	
	
\newcommand{\eexo}{\vspace{5mm}} % espace pour séparer les exercices


\begin{document}
%		\title{\vspace{-3cm}Série 1}
%		\date{\vspace{-2cm}}
%		\maketitle

\fancyhead[CO,CE]{\sc Série \arabic{section} \hspace{0.5mm}}

\setcounter{section}{\numero}
\section{Types et opérations}				
\Opensolutionfile{mycor}[cor_01]

\exo{}  ~\\ 
Tester dans la console \py les commandes suivantes
\begin{multicols}{3}
	\begin{lstlisting}[numbers=none]
>>> 2-9
>>> 7 + 3 * 4
>>> (7 + 3) * 4
>>> 20 / 3
>>> 20 // 3
>>> 18 / 3
>>> 18 // 3
>>> 20 % 3
>>> 18 % 3
>>> 4 * 4 * 4
>>> 4 ** 3
>>> 5 ** 1 * 2
>>> 5 ** (1 * 2)
>>> 5 * 1 ** 3
>>> (5 * 1) ** 3
>>> 59 * 100 // 60
>>> 59 * (100 // 60)
	\end{lstlisting}
\end{multicols}
\begin{enumerate}
	\item Que représentent les symboles \lstinline{**} et \lstinline{%} ?
		\item Pouvez-vous expliquer la différence entre \lstinline{/} et \lstinline{//} ?\\
	\end{enumerate}

	\begin{correction}
		~\\ 
		\begin{enumerate}
	\item Le symbole \lstinline{**} représente l'exponentiation (\lstinline{2**3} correspond à $2^3$) et \lstinline{%} le reste de la division entière.
	\item Le symbole \lstinline{/} représente la division décimale ($7/2$ donne 3.5 ) tandis que \lstinline{//} calcule la division entière ($7//2$ donne 3 ).
 		\end{enumerate}
	\end{correction}
\finexo

\exo{}  ~\\ 
Tester dans la console \py les commandes suivantes
\begin{lstlisting}[numbers=none]
>>> type(123)
>>> type(123.)
>>> type("blabla")
>>> type(True)
\end{lstlisting}

Indiquer lesquels des objets suivants sont valides en Python et le cas échéant de quel type d'objet il s'agit :
\begin{multicols}{2}
	\begin{enumerate}[label=\alph*)]
		\item "rewr"
		\item gdru
		\item 5
		\item 5.0
		\item 5,0
		\item 'sism2021'
		\item 18gd
		\item "65.5"
		\item True
		\item "False"
	\end{enumerate}
\end{multicols}
	\begin{correction}
	~\\ 
	Utiliser la fonction \lstinline{type()} dans l'interpréteur \py pour vérifier vos réponses.
\end{correction}
\finexo

\newpage

\exo{}[Convertir dans un type]  ~\\ 
 En s'aidant de la console \py, essayez de déterminer le rôle des fonctions \lstinline{int()},
\lstinline{float()} et \lstinline{str()} en écrivant par exemple les lignes suivantes :
\begin{lstlisting}[numbers=none]
>>> int(3.6)
>>> float(3)
>>> str(3)
>>> 2**1000
>>> float(2**1000)
>>> int('3')
\end{lstlisting}
\vspace{0.3cm}
\textit{Indication : Vous pouvez vous aider de la commande \lstinline{type()} pour analyser le résultat des fonctions \lstinline{int()}, \lstinline{float()} et \lstinline{str()}.}
	\begin{correction}
	~\\ 
	Les fonctions \lstinline{int()}, \lstinline{float()} et \lstinline{str()} permettent de convertir les données dans le type indiqué par la fonction.
\end{correction}
\finexo

\exo{}[Expression et priorité des opérations]  ~\\ 
Prédire le résultat des instructions suivantes en \py
\begin{multicols}{2}
	\begin{enumerate}[label=\alph*)]
		\item \lstinline{3+2*5}
		\item \lstinline{2*5**2}
		\item \lstinline{(2*5)**2}
		\item \lstinline{3*5+3}
		\item \lstinline{6*2/4}
		\item \lstinline{6/4*2}
		\item \lstinline{6.0/4*2}
		\item \lstinline{7 + 3 * 6 / 2 - 1}
		\item \lstinline{2 % 2 + 2 * 2 - 2 / 2}
			\item \lstinline{3 * 9 * (3 + (9 * 3 / (3))))}
		\end{enumerate}
	\end{multicols}
	Entrer ces expressions dans un terminal \py pour vérifier les réponses.
	
	\begin{correction}
	~\\ 
\begin{multicols}{2}
	\begin{enumerate}[label=\alph*)]
		\item \lstinline{13}
		\item \lstinline{50}
		\item \lstinline{100}
		\item \lstinline{18}
		\item \lstinline{3.0}
		\item \lstinline{3.0}
		\item \lstinline{3.0}
		\item \lstinline{15.0}
		\item \lstinline{3.0}
		\item \lstinline{324.0}
		\end{enumerate}
	\end{multicols}
\end{correction}
\finexo


% Solution		
		\newpage
		\setcounter{page}{1}
		\setcounter{section}{\numero}
		\Closesolutionfile{mycor}
		\titleformat{\section}[hang]{\Large \bfseries}{Corrigé Série \thesection:\ }{0pt}{}
		

		\fancyhead[CO,CE]{\sc Corrigé Série \arabic{section} \hspace{0.5mm}}
		\section{}
		\Readsolutionfile{mycor}
	\end{document}